% Generated draft: requires mathematical and source review.
\documentclass[11pt]{article}
\usepackage[margin=25mm]{geometry}
\usepackage{fontspec}
\setmainfont{Latin Modern Roman}
\usepackage{amsmath,amssymb,mathtools,unicode-math}
\setmathfont{Latin Modern Math}
\usepackage{enumitem,needspace,xurl,hyperref}
\setlength{\parindent}{0pt}
\setlength{\parskip}{5pt}
\setlength{\emergencystretch}{3em}
\newcommand{\N}{\mathbb{N}}
\newcommand{\Z}{\mathbb{Z}}
\newcommand{\Q}{\mathbb{Q}}
\newcommand{\R}{\mathbb{R}}
\newcommand{\C}{\mathbb{C}}
\newcommand{\F}{\mathbb{F}}
\newcommand{\A}{\mathbb{A}}
\newcommand{\PP}{\mathbb{P}}
\newcommand{\E}{\mathbb{E}}
\newcommand{\eps}{\varepsilon}
\newcommand{\dd}{\,\mathrm{d}}
\begin{document}
\section*{UnsolvedMath ID: 30003781}
\textbf{Source problem code:} catalogue entry
\subsection{Definitions and mathematical statement}
A closed manifold is a compact manifold without boundary; connectedness is not specified in the supplied record. An action \(\Gamma\curvearrowright M\) is free if \(\gamma x=x\) implies \(\gamma=e\), minimal if every orbit is dense in \(M\), and isometric if every group element preserves the metric on \(M\). Here finite presentation is understood as a condition on the acting groups \(\Gamma\) and \(\Lambda\).

Two groups are commensurable if they have isomorphic finite-index subgroups. The exclusion of a nontrivial free abelian factor is understood to mean that neither acting group is commensurable to any group of the form \(H\times\Z^k\), where \(H\) is a group and \(k\geq1\).

The notation \(M_t\) and \(N_s\) is not defined in the supplied statement. A usual warped-level interpretation is the following: given the metric \(d_M\) underlying the action and a finite symmetric generating set \(S_\Gamma\), equip \(M\) with the largest metric \(d_{M,t}\) satisfying
\[
 d_{M,t}(x,y)\leq t\,d_M(x,y),\qquad
 d_{M,t}(x,\gamma x)\leq1
 \quad(\gamma\in S_\Gamma).
\]
Define \(d_{N,s}\) analogously. This interpretation is supplied here to explain the warped-cone notation, but the record does not specify the metrics, their regularity, or the precise warped-level convention; these require confirmation against the cited source.

Under this interpretation, uniform quasi-isometries along sequences mean that there are positive numbers \(t_n,s_n\), maps \(f_n:M_{t_n}\to N_{s_n}\), and constants \(L\geq1\), \(C\geq0\), independent of \(n\), such that
\[
 t_n\longrightarrow\infty,\qquad s_n\longrightarrow\infty,
\]
\[
 L^{-1}d_{M,t_n}(x,y)-C
 \leq d_{N,s_n}(f_n(x),f_n(y))
 \leq Ld_{M,t_n}(x,y)+C
 \quad(x,y\in M),
\]
and every point of \(N_{s_n}\) is at distance at most \(C\) from \(f_n(M)\). No relation between the rates of growth of \(t_n\) and \(s_n\) is stated.

\emph{Source question, with unresolved terminology retained.} Let \(\Gamma\) and \(\Lambda\) be finitely presented groups acting freely, minimally, and isometrically on closed manifolds \(M\) and \(N\), respectively. Suppose that neither group is commensurable to a group having a nontrivial free abelian direct factor, and that \(\Gamma\) has Property FSL. If uniform quasi-isometries \(M_{t_n}\to N_{s_n}\) exist for sequences \(t_n,s_n\to\infty\), must the two actions admit a commensuration?

The supplied record does not define Property FSL or a commensuration of actions. In particular, group commensurability as defined above does not define the requested relation between actions. The question therefore remains a qualified transcription rather than a fully specified mathematical proposition; the intended definitions and any associated hypotheses must be recovered from the cited source before mathematical review.
\subsection{Short English statement}
The question asks whether uniform quasi-isometries between increasingly large warped levels force the underlying actions to admit a commensuration, under the stated freeness, minimality, finite-presentation, abelian-factor, and Property FSL assumptions. The intended meanings of Property FSL, commensuration of actions, and the warped-level metrics require source confirmation. The supplied open label is dated curation, not an independently verified conclusion.
\subsection{Sources}
\begin{itemize}
\item Local UnsolvedMath registry, \texttt{data/unsolvedmath/problems.json}, supplied complete record for permanent ID \(30003781\), titled \emph{Warped-Cone Quasi-Isometries and Commensuration of Actions}; exact source code \texttt{catalogue entry}. The record gives proposed year 2018 and names no proposer. Its labels \texttt{open} and \texttt{corrected\_verified}, including the claim that omitted definitions were repaired, are supplied registry assertions rather than independent verification here.
\item Supplied report citation: \emph{Mini-Workshop: Superexpanders and Their Coarse Geometry (2019)}, Oberwolfach Reports, pp.\ 1117--1160, DOI \texttt{10.4171/owr/2018/19}, \url{https://doi.org/10.4171/owr/2018/19}. The literature-source list calls this the OWR 2018 report. The differing year descriptions are retained without resolution. The supplied background describes Oberwolfach Reports as documenting workshops and their open-problem discussions.
\item Supplied literature source: Fisher--Nguyen--van Limbeek, warped cones, \url{https://doi.org/10.1016/j.aim.2019.02.015}.
\item Supplied literature source: Vigolo, warped cones, \url{https://arxiv.org/abs/1710.04955}.
\item Literature assessment dated exactly \texttt{2026-08-21}: the supplied curation labels the problem open and reports that Fisher--Nguyen--van Limbeek prove strong rigidity for free minimal actions without abelian factors, while later warped-cone work leaves the broad relationship between coarse geometry and dynamics unresolved. This is supplied dated curation derived from the OWR collection, not an independently verified literature assessment. Its applicability to the exact question and its time-sensitive claims require independent verification.
\end{itemize}
\end{document}
